\subsection{Опорные векторы}

\begin{definition}\label{def:support-vector}
Допустимая точка $x_0[N]$ называется \textbf{опорным вектором} (\ov), если столбцы
$\{a_j[M]\}_{j \in \Nzp}$, отвечающие положительным компонентам, линейно независимы.
\end{definition}

\begin{remark}
Из $\rank A[M,N] = \abs{M}$ следует $\abs{\Nzp} \le \abs{M}$: линейно независимых
столбцов в $\R^{m}$ не может быть больше $m$.
\end{remark}

\begin{definition}\label{def:degenerate}
Опорный вектор $x_0[N]$ называется \textbf{невырожденным}, если $\abs{\Nzp} = \abs{M}$,
и \textbf{вырожденным} иначе (то есть при $\abs{\Nzp} < \abs{M}$).
\end{definition}

Отсюда два случая.

\begin{enumerate}[label=(\arabic*),left=0pt]
  \item Если \ov{} невырожден, то $\det A[M, \Nzp] \ne 0$ для этого вектора:
        квадратная матрица $\abs{M} \times \abs{M}$ из линейно независимых столбцов
        невырождена. Здесь $\Np := \Nzp$.
  \item Если \ov{} вырожден, то набор $\Nzp$ можно \emph{пополнить} до множества $\Np$
        с $\abs{\Np} = \abs{M}$ так, что $\det A[M, \Np] \ne 0$.
\end{enumerate}

В терминах компонент самого вектора это даёт
\[
  (1) \implies x_0[i] > 0 \quad \forall i \in \Nzp = \Np ;
\]
\[
  (2) \implies
  \begin{cases}
    x_0[i] > 0, & i \in \Nzp, \\
    x_0[i] = 0, & i \in \Np \setminus \Nzp .
  \end{cases}
\]

\begin{idea}
Пополнение в случае (2) возможно именно из-за условия $\rank A[M,N] = \abs{M}$:
линейно независимый набор $\{a_j\}_{j \in \Nzp}$ можно достроить до базиса
пространства $\R^{m}$ столбцами той же матрицы $A$ (теорема о дополнении до базиса).

Разница между случаями — в том, сколько нулевых компонент попало в базис. У
невырожденного \ov{} базис состоит ровно из столбцов с положительными компонентами;
у вырожденного в базис добавляются «лишние» столбцы, которым отвечают нули. Поэтому
одному вырожденному \ov{} может отвечать \emph{несколько} разных базисов — это будет
видно в примере~\ref{ex:six-bases}.
\end{idea}

\begin{definition}\label{def:basis}
Набор столбцов $\{a_j[M]\}_{j \in \Np}$ называется \textbf{базисом опорного вектора}.
\end{definition}

\begin{lemma}[о числе опорных векторов]\label{lem:count}
Число опорных векторов конечно и не превосходит
\[
  \binom{n}{m} = \frac{n!}{m!\,(n-m)!} .
\]
\end{lemma}

\inputfig{figures/fig-matrix-selection}

\begin{proof}
Выделим из $A[M,N]$ всевозможные квадратные подматрицы вида $A[M,M']$,
составленные из $\abs{M} = m$ столбцов. Если $\det A[M,M'] = 0$, такая подматрица
исключается из рассмотрения; остаются только подматрицы с ненулевым определителем.
Всего наборов из $m$ столбцов, выбранных из $n$, ровно $\binom{n}{m}$, значит
невырожденных подматриц не больше $\binom{n}{m}$.

Для каждой оставшейся подматрицы решаем систему
\[
  A[M,M']\, x[M'] = b[M] ,
\]
полагая остальные компоненты нулями. Решение единственно, поскольку определитель
$\det A[M,M']$ отличен от нуля. Если полученный вектор удовлетворяет ограничению на знак,
$x[i] \ge 0$ для всех $i \in M'$, то он принимается к рассмотрению. Так получаются
все опорные векторы, и их не больше числа перебранных базисов.
\end{proof}

\begin{remark}
Одному взятому базису, составленному из столбцов $A$, соответствует \emph{ровно один}
опорный вектор. Обратное неверно: вырожденному \ov{} может отвечать несколько
базисов, поэтому оценка $\binom{n}{m}$ не точна.
\end{remark}

\begin{example}\label{ex:six-bases}
Рассмотрим задачу
\[
  \begin{cases}
    5x_1 + 5x_2 + 6x_3 + 4x_4 = 5, \\
    4x_1 + 5x_2 - 5x_3 - 4x_4 = 4, \\
    x_i \ge 0, \quad i = \overline{1,4} .
  \end{cases}
\]
Здесь $m = 2$, $n = 4$, столбцы матрицы
\[
  a_1 = \begin{pmatrix} 5 \\ 4 \end{pmatrix}, \quad
  a_2 = \begin{pmatrix} 5 \\ 5 \end{pmatrix}, \quad
  a_3 = \begin{pmatrix} 6 \\ -5 \end{pmatrix}, \quad
  a_4 = \begin{pmatrix} 4 \\ -4 \end{pmatrix} ,
\]
и по лемме~\ref{lem:count} перебираем не более $\binom{4}{2} = 6$ базисов.

\medskip
\noindent\textbf{1)} Базис $a_1, a_2$ (остальные переменные — нули):
\[
  \begin{cases} 5x_1 + 5x_2 = 5 \\ 4x_1 + 5x_2 = 4 \end{cases}
  \implies (x_1, x_2) = (1, 0),
  \qquad x_{\mathrm{I}}\T = (1, 0, 0, 0) \ \text{— \ov} .
\]

\noindent\textbf{2)} Базис $a_1, a_3$:
\[
  \begin{cases} 5x_1 + 6x_3 = 5 \\ 4x_1 - 5x_3 = 4 \end{cases}
  \implies (x_1, x_3) = (1, 0),
  \qquad x_{\mathrm{II}}\T = (1, 0, 0, 0) \ \text{— \ov} .
\]

\noindent\textbf{3)} Базис $a_1, a_4$:
\[
  \begin{cases} 5x_1 + 4x_4 = 5 \\ 4x_1 - 4x_4 = 4 \end{cases}
  \implies (x_1, x_4) = (1, 0),
  \qquad x_{\mathrm{III}}\T = (1, 0, 0, 0) \ \text{— \ov} .
\]

\noindent\textbf{4)} Базис $a_2, a_3$:
\[
  \begin{cases} 5x_2 + 6x_3 = 5 \\ 5x_2 - 5x_3 = 4 \end{cases}
  \implies (x_2, x_3) = \Big(\tfrac{49}{55},\ \tfrac{1}{11}\Big),
  \qquad x_{\mathrm{IV}}\T = \Big(0, \tfrac{49}{55}, \tfrac{1}{11}, 0\Big)
  \ \text{— \ov} .
\]

\noindent\textbf{5)} Базис $a_2, a_4$:
\[
  \begin{cases} 5x_2 + 4x_4 = 5 \\ 5x_2 - 4x_4 = 4 \end{cases}
  \implies (x_2, x_4) = \Big(\tfrac{9}{10},\ \tfrac{1}{8}\Big),
  \qquad x_{\mathrm{V}}\T = \Big(0, \tfrac{9}{10}, 0, \tfrac{1}{8}\Big)
  \ \text{— \ov} .
\]

\noindent\textbf{6)} Базис $a_3, a_4$:
\[
  \begin{cases} 6x_3 + 4x_4 = 5 \\ -5x_3 - 4x_4 = 4 \end{cases}
  \implies (x_3, x_4) = \Big(9,\ -\tfrac{49}{4}\Big),
  \qquad x_{\mathrm{VI}} \ \text{не существует, т.\,к. } -\tfrac{49}{4} < 0 .
\]
\end{example}

\begin{center}
\begin{tabular}{clll}
\toprule
№ & базис & решение & опорный вектор \\
\midrule
1 & $a_1, a_2$ & $(x_1, x_2) = (1, 0)$ & $(1,\,0,\,0,\,0)$ \\
2 & $a_1, a_3$ & $(x_1, x_3) = (1, 0)$ & $(1,\,0,\,0,\,0)$ \\
3 & $a_1, a_4$ & $(x_1, x_4) = (1, 0)$ & $(1,\,0,\,0,\,0)$ \\
4 & $a_2, a_3$ & $(x_2, x_3) = \big(\tfrac{49}{55}, \tfrac{1}{11}\big)$
  & $\big(0,\,\tfrac{49}{55},\,\tfrac{1}{11},\,0\big)$ \\
5 & $a_2, a_4$ & $(x_2, x_4) = \big(\tfrac{9}{10}, \tfrac{1}{8}\big)$
  & $\big(0,\,\tfrac{9}{10},\,0,\,\tfrac{1}{8}\big)$ \\
6 & $a_3, a_4$ & $(x_3, x_4) = \big(9, -\tfrac{49}{4}\big)$ & нет: $x_4 < 0$ \\
\bottomrule
\end{tabular}
\end{center}

\begin{idea}
Из шести базисов годными оказались пять, но \emph{различных} опорных векторов только
три:
\[
  (1,0,0,0), \qquad \Big(0, \tfrac{49}{55}, \tfrac{1}{11}, 0\Big),
  \qquad \Big(0, \tfrac{9}{10}, 0, \tfrac{1}{8}\Big) .
\]
Первый из них вырожден: $\abs{\Nzp} = 1 < 2 = \abs{M}$, и ему отвечают сразу три
базиса ($a_1,a_2$; $a_1,a_3$; $a_1,a_4$) — все они содержат обязательный столбец
$a_1$ и один «лишний», которому соответствует нулевая компонента. Два других \ov{}
невырождены, у каждого ровно один базис. Это и есть обещанная неточность оценки
$\binom{n}{m}$: пять годных базисов против трёх опорных векторов.
\end{idea}
